\parthead{VII}{What the landing requires}\label{part:7}
The 47-degree mismatch is easier to understand after following the motion to its end. A modest bank accumulates a substantial turn because lift acts over time. Radius and turning rate set the space and time needed. Descent and flare must then deliver a touchdown state from which the remaining runway can absorb the forward motion.

The two constructed landings put these relationships into one calculation. Each joins an approach, a flare and a full stop, with a separate continuous-time certificate. They expose choices hidden by a simple drawing: the synthetic entry and corridor, the allowance for a small rise, the scaling of competing constraints and the stipulated braking law. An explicit path allows those choices to be examined individually.

Within the restricted calm-air family, compactness and regularity also establish that a best primary member exists, together with a secondary effort minimum on the retained-primary set. The proof in Appendix~\ref{app:bounds} establishes attainment without identifying the optimizers. The remaining gaps between achieved scores and necessary upper bounds concern this mathematical family; they do not establish that the historical procedure was suboptimal.

\subsection*{Putting the calculation back into the city}
The historical images make the next challenge tangible. The recovered surface has roof forms, roads and openings, but its vertical origin remains local and much of its geometry lies outside the control hull. A clearance calculation needs a common runway/obstacle datum, dated terrain and roof coverage, physical upper bounds, and the aircraft's swept volume with a defensible tracking allowance. A more complete rendering helps only when additional geometry comes from additional evidence.

The aircraft and braking model must also be calibrated to the intended configuration and conditions, then tested against routine flights. Uncertain wind requires a causal control policy using information available at the time; choosing a different curve after each wind is known does not supply one. Go-around viability and physically supported outcome uncertainty remain part of that larger task. Appendix~\ref{app:scope} gives the full research requirements and related methods, while Appendix~\ref{app:ground} derives the conditional stopping bounds.

The opening question was how a large aircraft could negotiate the turn and still have enough space, time and runway to land. The calculations explain how those demands fit together. Bank redirects lift; sustained turning accumulates direction; descent changes the energy balance; wind separates air motion from ground motion; touchdown commits the remaining pavement to delay and braking. Each stage constrains the next.

Under the declared constraints, turn, descent and stop fit within one continuous mathematical construction. Connecting that construction to a historical landing requires the dated city and a calibrated aircraft. Kai Tak drew the eye to the turn. The calculation keeps following the aircraft until it stops.
